%%
%% This is file `acmtog.tex',
%% generated with the docstrip utility.
%%
%% The original source files were:
%%
%% samples.dtx  (with options: `all,acmtog')
%%
%% IMPORTANT NOTICE:
%%
%% For the copyright see the source file.
%%
%% Any modified versions of this file must be renamed
%% with new filenames distinct from acmtog.tex.
%%
%% For distribution of the original source see the terms
%% for copying and modification in the file samples.dtx.
%%
%% This generated file may be distributed as long as the
%% original source files, as listed above, are part of the
%% same distribution. (The sources need not necessarily be
%% in the same archive or directory.)
%%
%%
%% Commands for TeXCount
%TC:macro \cite [option:text,text]
%TC:macro \citep [option:text,text]
%TC:macro \citet [option:text,text]
%TC:envir table 0 1
%TC:envir table* 0 1
%TC:envir tabular [ignore] word
%TC:envir displaymath 0 word
%TC:envir math 0 word
%TC:envir comment 0 0
%%
%% The first command in your LaTeX source must be the \documentclass
%% command.
%%
%% For submission and review of your manuscript please change the
%% command to \documentclass[manuscript, screen, review]{acmart}.
%%
%% When submitting camera ready or to TAPS, please change the command
%% to \documentclass[sigconf]{acmart} or whichever template is required
%% for your publication.
%%
%%
\documentclass[acmtog,nonacm]{acmart}   % nonacm: no ACM reference box, no journal footer
%%
%% \BibTeX command to typeset BibTeX logo in the docs
\AtBeginDocument{%
  \providecommand\BibTeX{{%
    Bib\TeX}}}

%% Rights management information.
%% No ACM boilerplate: no copyright/permission block, no "Authors' Contact Information" footnote
\setcopyright{none}
\authorsaddresses{}
% \copyrightyear{2026}
% \acmYear{2026}
% \acmDOI{XXXXXXX.XXXXXXX}

\citestyle{acmauthoryear}
\usepackage{booktabs}

\begin{document}

\title{FiberCANDI - Deconvolving Bulk Assays to Single Molecule Resolution}

\author{Azaan Rehman}
\email{ara199@sfu.ca}
\correspondingauthor
\affiliation{%
  \institution{Simon Fraser University}
  \city{Burnaby}
  \state{British Columbia}
  \country{Canada}
}

\author{Mehdi Foroozandeh}
\affiliation{%
  \institution{Simon Fraser University}
  \city{Burnaby}
  \state{British Columbia}
  \country{Canada}
}

\author{Andrew B. Stergachis}
\affiliation{%
  \institution{University of Washington}
  \city{Seattle}
  \state{Washington}
  \country{USA}
}

\author{Christine Queitsch}
\affiliation{%
  \institution{University of Washington}
  \city{Seattle}
  \state{Washington}
  \country{USA}
}

\author{Maxwell W. Libbrecht}
\affiliation{%
  \institution{Simon Fraser University}
  \city{Burnaby}
  \state{British Columbia}
  \country{Canada}
}

\renewcommand{\shortauthors}{Rehman et al.}

\maketitle

\section{Research problem}

\textbf{Question.} Can we impute histone modifications and CTCF binding on \emph{individual} DNA molecules from their single-molecule chromatin accessibility and DNA sequence, when the only supervision available is population-averaged (bulk) data? And can such imputations be verified without molecule-level ground truth?

\textbf{Background for non-biologists.} (Jump to computational problem to come to more CS terms)
\begin{itemize}
  \item \emph{Chromatin and nucleosomes.} In a cell, DNA is wrapped around protein spools called \emph{nucleosomes}, about 147 bases each, like beads on a string. The DNA plus its proteins is \emph{chromatin}. A \emph{fiber} is one such DNA molecule, which in our case is a few thousand to tens of thousands of bases long.
  \item \emph{Accessibility.} Stretches of DNA between nucleosomes are \emph{accessible}: proteins can reach them. Regulatory elements (switches that turn genes on or off) are usually accessible in the cells where they are active. ATAC-seq measures accessibility in bulk.
  \item \emph{Histone marks.} Nucleosomes carry chemical tags (e.g.\ H3K27ac, H3K4me3, H3K27me3), named after the protein (histone H3), the position (lysine 27) and the tag (acetyl, methyl). Different marks indicate active promoters, active enhancers, transcribed genes or silenced regions. \emph{CTCF} is a DNA-binding protein that organizes the genome into loops. Although the cells share the same genome, the histone marks vary significantly for different types of cells.
  \item \emph{Bulk assays.} ChIP-seq and ATAC-seq pool DNA from millions of cells and report, at each position, how enriched a mark is relative to a control (\emph{fold change}). The result is a \emph{track}: one number per base, averaged over all molecules. ENCODE provides such tracks for many cell lines (cultured cell populations such as GM12878, a B-cell line, and HepG2, a liver line).
  \item \emph{Fiber-seq.} A new assay that uses an enzyme that adds a chemical tag (m6A) to accessible bases. Long-read sequencing then reads each molecule end to end, along with its tags, so accessibility, nucleosome positions and DNA methylation (at CpG sites) are known for \emph{every individual molecule}, not just on average.
\end{itemize}

\textbf{Motivation.} Fiber-seq \citep{Stergachis20} reads multi-kilobase chromatin fibers and records accessibility (m6A), CpG methylation and nucleosome positions on each molecule, with haplotype phase. Histone marks, by contrast, are measured almost only in bulk (ChIP-seq). Bulk averages over millions of molecules and hides heterogeneity: allele-specific states, sub-populations, and which marks co-occur on one molecule. Deconvolving bulk assays onto single molecules would turn existing Fiber-seq data into a multi-mark single-molecule atlas.

\textbf{Computational problem.} The overview is as such: We start with ternary three-dimensional tensor and predict a real valued three-dimensional vector. We don't have any ground truth on this resolution so we average over one dimension to get a two-dimensional vector where we do have the ground truth, and from there we derive the loss and backpropagate in a weakly-supervised fashion. For a genomic window of $L=5$\,kb, the input is a stack of $N\le100$ fibers, $X\in\{-1,0,1\}^{C\times L\times N}$, with $C=5$ channels (m6A, CpG, accessible patch, nucleosome, FIRE element; called with fibertools \citep{Jha24}), plus the one-hot reference sequence $S\in\{0,1\}^{4\times L}$. For each fiber $f$ and assay $k$ ($K=8$: ATAC, CTCF and six histone marks), the model outputs a per-base value $z_{f,k}\in\mathbb{R}_{\ge0}^{L}$. Fibers are processed independently, and only their average is supervised, against the ENCODE fold-change (FC) track $y_k\in\mathbb{R}_{\ge0}^{L}$ \citep{ENCODE20} which acts as the ground truth; a real valued one dimensional signal over the genome. Bulk read counts add up over molecules, so the average is taken in FC space:
\begin{equation}
\hat y_k(x) = \operatorname{asinh}\Big(\tfrac{1}{|F_x|}\textstyle\sum_{f\in F_x}\sinh z_{f,k}(x)\Big),
\end{equation}
\begin{equation}
\mathcal{L}=\textstyle\sum_{k,x} w_{k,x}\,\big(\hat y_k(x)-\operatorname{asinh} y_k(x)\big)^2 .
\end{equation}
Here $F_x$ is the set of fibers covering position $x$, each trimmed to its first through last nucleosome, and $w$
restricts the loss to positions covered by at least 5 fibers. The per-fiber outputs $z_{f,k}$ are the imputation results.

\textbf{Assumptions.} (A1) A molecule's accessibility and nucleosome pattern, together with its sequence, carry information about its marks. (A2) A FC track is linear in the fraction of marked molecules, up to a per-experiment scale. (A3) The Fiber-seq and ENCODE samples of a cell line describe the same population.

\textbf{Evidence plan.} Our hypothesis (H) is that the per-fiber outputs carry information specific to each molecule. A model can also fit bulk tracks by predicting the bulk on every molecule (our null, H0), so bulk accuracy is necessary but not sufficient. Throughout, $p$ is a predicted bulk track (the average of a cell type's fibers, $\hat y_k$ in Eq.~1), $t$ is that cell type's ENCODE track ($\operatorname{asinh} y_k$), and all metrics pool every scored base of the test chromosome. We test H against H0 at three levels.

(i) \emph{Bulk recovery} on held-out chromosomes and cell types. We use 5 development and 3 held-out cell lines; the held-out ones range from near to far from the development set. We report the Pearson correlation $r_{pt}\in[-1,1]$ (1: the predicted shape along the genome matches the target up to scale and offset; 0: no linear relation; $-1$: inverted) and the mean squared error between $p$ and $t$ ($\ge0$, in squared asinh units; 0 is a perfect prediction). As a baseline we use the \emph{cross-cell average} $b$: at each position, the mean of the training cell lines' tracks, which needs no information about the evaluated cell type. Much of every track is shared across cell types, so $b$ is a strong baseline. To isolate what is specific to the cell type, we report the \emph{cell-specific correlation}, the partial correlation of $p$ and $t$ given $b$. Both are regressed on $b$, and their residuals are correlated: $r(p,t\mid b)=\operatorname{corr}\big(p-\hat p(b),\,t-\hat t(b)\big)=(r_{pt}-r_{bt}r_{pb})/\sqrt{(1-r_{bt}^2)(1-r_{pb}^2)}$. It lies in $[-1,1]$: 1 means the model predicts exactly where and how this cell type departs from the average; 0 means nothing beyond the average (the baseline itself, or any rescaling of it, scores exactly 0); negative values mean the predicted departures point the wrong way.

(ii) \emph{Mixture-only supervision} (main test). We train only on in-silico mixtures of two cell types $A$ and $B$ with varying proportions. Each training sample pools both types' fibers at one locus, and its target is the mixture's track; neither type's own track is ever used. We then run each cell type's fibers alone, giving $p_A$ and $p_B$, and compare them with $t_A$ and $t_B$. Under H0 the model predicts the mixture for both, so $p_A\approx p_B$. Each readout has a known value under H0:
\begin{itemize}
  \item \emph{Differential correlation} $\operatorname{corr}(p_A-p_B,\,t_A-t_B)$, per base over positions scored in both cell types: does the predicted difference between the two cell types follow the true one? Range $[-1,1]$: 1 means the predicted differences follow the true ones exactly (up to scale), 0 means no relation (value for H0), and $-1$ means they are systematically reversed.
  \item \emph{Sign accuracy}: over \emph{differential bins} (1\,kb bins where $|t_A-t_B|\ge0.5$), the fraction where the mean of $p_A-p_B$ over the bin has the same sign as $t_A-t_B$. Range $[0,1]$: 1 means the higher cell type is identified at every bin, 0.5 is chance (value under H0), and 0 means it is always reversed.
  \item \emph{Fiber-origin AUC}: at each differential bin, score each fiber by its mean output $z_f$ over the bin. The AUC is the probability that a random fiber from the cell type with the higher track scores above a random fiber from the other; we average it over bins. It asks whether a \emph{single molecule's} prediction reveals where it came from. Range $[0,1]$: 1 means every fiber of the cell type with the mark outranks every fiber of the other (perfect separation), 0.5 means no separation (value under H0), and 0 means perfect separation in the wrong direction. Even a perfect model stays below 1, because not every molecule of a cell type carries its mark.
\end{itemize}
As a control for library effects, we use two Fiber-seq libraries of one cell line, which share one true track: any separation between them comes from library differences, not biology.

(iii) \emph{Within-population structure}, from held-out CpG methylation and X inactivation, where molecules of a single cell line must differ.

\textbf{Preliminary results.} Over all 6,843 test tiles of chr21, the cell-specific correlation is positive in all 8 cell lines ($+0.19$ to $+0.41$; $+0.19$ to $+0.27$ in cell lines never seen in training). In the mixture test (GM12878 + HepG2), the model trained only on mixtures separates the two populations' molecules as well as a model trained on each cell type's own track (Table~\ref{tab:t3}). The separation is locus-specific: it reverses with the direction of the difference. For two libraries of one cell line, $|\mathrm{AUC}-0.5|=0.04$, against $0.18$ for GM12878 vs HepG2.

\begin{table}[h]
\caption{Mixture test, mean over 8 assays.}
\label{tab:t3}
\small
\begin{tabular}{lccc}
\toprule
 & H0 & Mixture only & Own tracks \\
\midrule
Differential $r$ (5\,kb windows) & 0 & 0.49 & 0.47 \\
Sign accuracy (differential bins) & 0.50 & 0.78 & 0.74 \\
Fiber-origin AUC & 0.50 & 0.64 & 0.60 \\
\bottomrule
\end{tabular}
\end{table}

\section{Literature review}

\textbf{Single-molecule chromatin profiling.} Fiber-seq \citep{Stergachis20} marks accessible DNA with a methyltransferase and reads it with long reads, giving accessibility for every molecule. DiMeLo-seq \citep{Altemose22} maps one protein or histone mark per experiment directly on single molecules. Single-cell CUT\&Tag \citep{Bartosovic21} profiles marks in single cells, but with sparse coverage and only one or a few marks per cell. Together, these show that \emph{measuring} many marks on one molecule is still out of reach. We instead \emph{impute} them, and DiMeLo-seq data are the natural ground truth for one mark at a time.

\textbf{Epigenome imputation.} ChromImpute \citep{Ernst15}, PREDICTD \citep{Durham18} and Avocado \citep{Schreiber20Avocado} predict missing bulk tracks from other tracks of the same cell type. They use regression or tensor factorization over cell types $\times$ assays $\times$ positions. CANDI \citep{Foroozandeh25} imputes and denoises bulk assays with self-supervision and calibrated confidence. All of them output population averages and need existing tracks of the target cell type. FiberCANDI keeps the target (ENCODE assays) but changes the unit of prediction from a track to a molecule, and needs only the target cell type's fibers. \citet{Schreiber20Pitfall} showed that models predicting across cell types often fail to beat the average across training cell types. We adopt that baseline and score the partial correlation given it.

\textbf{Sequence-to-function models.} Enformer \citep{Avsec21} predicts thousands of tracks from DNA alone, using long-range attention. Its cell-type specificity comes from one output head per training cell type, so it cannot describe a new cell type or individual molecules. In our model the DNA is shared context, and the fibers carry everything specific to the cell and the molecule. A pretrained DNA encoder is a planned extension.

\textbf{Weak supervision.} Learning from label proportions \citep{Quadrianto09} recovers labels of individual instances from bag-level aggregates, which is possible only when bags differ in composition. Our setting is a regression analogue: a bag is a fiber stack and its label is the bulk track. This is why the mixture test varies the proportions: varying them makes each population's signal identifiable.

Compared with these works, we connect bulk-assay imputation (the second theme) to single-molecule inputs (the first), and we give a falsifiable verification protocol for a setting with no molecule-level labels.

\section{Venue, contributions, and bibliography}

\textbf{Venue.} MLCB 2027 (Machine Learning in Computational Biology). Its focus on new machine-learning formulations with biological evaluation matches our contribution: a weakly supervised model and a falsifiable test of what the model learns.

\textbf{Prior work.} Two lines of prior work underlie this project. Bulk-assay imputation with CANDI \citep{Foroozandeh25} defines the assay-imputation setting we extend to single molecules. Fiber-seq \citep{Stergachis20} provides the single-molecule chromatin data and its processing \citep{Jha24}. Bulk targets and cCRE annotations come from ENCODE \citep{ENCODE20}.

\textbf{Contributions.} The work in this proposal is mine (A.R.):
\begin{enumerate}
  \item The problem formulation, and a reproducible data pipeline: ENCODE track selection; train, validation and test regions on separate chromosomes that avoid blacklisted regions and assembly gaps; and a principled split of cell types into 5 for development and 3 held out.
  \item Fiber aggregation consistent with fold change (Eq.~1). It gives each per-molecule output a physical meaning and makes mixtures exact.
  \item An evaluation framework that separates bulk recovery from single-molecule information: the cross-cell baseline with partial correlation, cell-specific and pairwise loci, and the mixture-only test with a library control.
  \item First evidence that the per-molecule imputations carry molecule-specific signal.
\end{enumerate}
To our knowledge, this is the first method to impute several histone marks per single molecule. Unlike CANDI and earlier imputation methods, its output is one track per molecule rather than one per cell type.

\bibliographystyle{ACM-Reference-Format}
\bibliography{proposal}

\end{document}

\endinput
